% SPDX-License-Identifier: CC-BY-SA-4.0
% Session: Form parallelism to concurrency
% Exercise: Amdahl's law

\begingroup

\begin{exercice}[Pile wait-free]
  On rappelle l'algorithme de pile wait-free inspiré de l'algorithme de Afek, Gafni et Morrison. 
  
  \begin{lstlisting}[numbers=none]
    public class WaitFreeStack<T> implements Stack<T> {

      private class Node {
        final AtomicReference<T> value = new AtomicReference<>();
        final int rank;
        final Node next;
        Node(Node next, int rank) {
          this.rank = rank;
          this.next = next;
        }
      }

      private final AtomicInteger last_index = new AtomicInteger(0);
      private final AtomicReference<Node> head = new AtomicReference<>(new Node(null, 0));

      public void push(T value) {
        if(value == null) throw new NullPointerException();
        int index = last_index.getAndIncrement();
        Node node = head.get();
        while(node.rank <= index) {
          head.compareAndSet(node, new Node(node, node.rank+1));
          node = head.get();
        }
        while(node.rank > index) node = node.next;
        node.value.set(value);
      }

      public T pop() {
        for(Node node = head.get(); node != null; node = node.next ) {
          T value = node.value.getAndSet(null);
          if(value != null) return value;
        }
        return null;
      }
      
    }
  \end{lstlisting}
  \begin{question}
  \item Modifiez l'algorithme pour supprimer les n\oe uds de la liste chaînée correspondant aux éléments qui ont déjà été dépilés. 
  \end{question}
\end{exercice}

\endgroup
\endinput
